% Generated by scripts/latex-perf-book.mjs. Do not edit by hand.
\section{Chapter 1: Baseline}
\label{sec:ch01}

Performance probe paragraph B01. This is included chapter 1 of twelve in
the large MDraft LuaLaTeX benchmark. Ordinary prose keeps the page count
realistic and gives the probe a stable word to edit. Equation~\ref{eq:ch12:integral} and
Table~\ref{tab:ch12:runs} are valid cross-chapter references.

\subsection{An integral}
The integral in Equation~\ref{eq:ch01:integral} is valid for every integer
$n \ge 0$. It is deliberately simple; the benchmark measures the engine, not
the mathematics.

\begin{equation}
  \label{eq:ch01:integral}
  \int_{0}^{1} x^{n}\,dx = \frac{1}{n+1}.
\end{equation}

The geometric series in Equation~\ref{eq:ch01:series} converges for
$|r| < 1$ and gives a second display to render in this chapter.

\begin{equation}
  \label{eq:ch01:series}
  \sum_{k=0}^{m} r^{k} = \frac{1-r^{m+1}}{1-r}.
\end{equation}

\subsection{A table}
\begin{table}[h]
  \centering
  \begin{tabular}{lrr}
    \toprule
    Run & Time & Pages \\
    \midrule
    One   & 1.20 & 3 \\
    Two   & 1.15 & 3 \\
    Three & 1.18 & 3 \\
    \bottomrule
  \end{tabular}
  \caption{A timing sample for chapter 1.}
  \label{tab:ch01:runs}
\end{table}

Table~\ref{tab:ch01:runs} is a valid float with a caption and a label. It
exists so the page builder has a boxed object that cannot be split across a
page break. The bibliography entry \cite{knuth1984} is resolved by the same compile.

\subsection{More prose}
A benchmark chapter needs enough text that the type\-set is visible above the
engine boot. This paragraph, the next one and the table above add up to a few
pages per chapter.

The content is intentionally repetitive across chapters. The point is not the
prose; it is that the include graph, the equation numbering and the table
labels all have to stay consistent while the probe edits one letter.

The final paragraph of chapter 1 gives the probe a second stable offset. A
change here must not invalidate the labels in any other chapter.

\subsection{Part 2}

This is included chapter 1 of twelve in
the large MDraft LuaLaTeX benchmark. Ordinary prose keeps the page count
realistic and gives the probe a stable word to edit. Equation~\ref{eq:ch12:k2:integral} and
Table~\ref{tab:ch12:k2:runs} are valid cross-chapter references.

\subsubsection{An integral}
The integral in Equation~\ref{eq:ch01:k2:integral} is valid for every integer
$n \ge 0$. It is deliberately simple; the benchmark measures the engine, not
the mathematics.

\begin{equation}
  \label{eq:ch01:k2:integral}
  \int_{0}^{1} x^{n}\,dx = \frac{1}{n+1}.
\end{equation}

The geometric series in Equation~\ref{eq:ch01:k2:series} converges for
$|r| < 1$ and gives a second display to render in this chapter.

\begin{equation}
  \label{eq:ch01:k2:series}
  \sum_{k=0}^{m} r^{k} = \frac{1-r^{m+1}}{1-r}.
\end{equation}

\subsubsection{A table}
\begin{table}[h]
  \centering
  \begin{tabular}{lrr}
    \toprule
    Run & Time & Pages \\
    \midrule
    One   & 1.20 & 3 \\
    Two   & 1.15 & 3 \\
    Three & 1.18 & 3 \\
    \bottomrule
  \end{tabular}
  \caption{A timing sample for chapter 1.}
  \label{tab:ch01:k2:runs}
\end{table}

Table~\ref{tab:ch01:k2:runs} is a valid float with a caption and a label. It
exists so the page builder has a boxed object that cannot be split across a
page break. The bibliography entry \cite{knuth1984} is resolved by the same compile.

\subsubsection{More prose}
A benchmark chapter needs enough text that the type\-set is visible above the
engine boot. This paragraph, the next one and the table above add up to a few
pages per chapter.

The content is intentionally repetitive across chapters. The point is not the
prose; it is that the include graph, the equation numbering and the table
labels all have to stay consistent while the probe edits one letter.

The final paragraph of chapter 1 gives the probe a second stable offset. A
change here must not invalidate the labels in any other chapter.

\subsection{Part 3}

This is included chapter 1 of twelve in
the large MDraft LuaLaTeX benchmark. Ordinary prose keeps the page count
realistic and gives the probe a stable word to edit. Equation~\ref{eq:ch12:k3:integral} and
Table~\ref{tab:ch12:k3:runs} are valid cross-chapter references.

\subsubsection{An integral}
The integral in Equation~\ref{eq:ch01:k3:integral} is valid for every integer
$n \ge 0$. It is deliberately simple; the benchmark measures the engine, not
the mathematics.

\begin{equation}
  \label{eq:ch01:k3:integral}
  \int_{0}^{1} x^{n}\,dx = \frac{1}{n+1}.
\end{equation}

The geometric series in Equation~\ref{eq:ch01:k3:series} converges for
$|r| < 1$ and gives a second display to render in this chapter.

\begin{equation}
  \label{eq:ch01:k3:series}
  \sum_{k=0}^{m} r^{k} = \frac{1-r^{m+1}}{1-r}.
\end{equation}

\subsubsection{A table}
\begin{table}[h]
  \centering
  \begin{tabular}{lrr}
    \toprule
    Run & Time & Pages \\
    \midrule
    One   & 1.20 & 3 \\
    Two   & 1.15 & 3 \\
    Three & 1.18 & 3 \\
    \bottomrule
  \end{tabular}
  \caption{A timing sample for chapter 1.}
  \label{tab:ch01:k3:runs}
\end{table}

Table~\ref{tab:ch01:k3:runs} is a valid float with a caption and a label. It
exists so the page builder has a boxed object that cannot be split across a
page break. The bibliography entry \cite{knuth1984} is resolved by the same compile.

\subsubsection{More prose}
A benchmark chapter needs enough text that the type\-set is visible above the
engine boot. This paragraph, the next one and the table above add up to a few
pages per chapter.

The content is intentionally repetitive across chapters. The point is not the
prose; it is that the include graph, the equation numbering and the table
labels all have to stay consistent while the probe edits one letter.

The final paragraph of chapter 1 gives the probe a second stable offset. A
change here must not invalidate the labels in any other chapter.

\subsection{Part 4}

This is included chapter 1 of twelve in
the large MDraft LuaLaTeX benchmark. Ordinary prose keeps the page count
realistic and gives the probe a stable word to edit. Equation~\ref{eq:ch12:k4:integral} and
Table~\ref{tab:ch12:k4:runs} are valid cross-chapter references.

\subsubsection{An integral}
The integral in Equation~\ref{eq:ch01:k4:integral} is valid for every integer
$n \ge 0$. It is deliberately simple; the benchmark measures the engine, not
the mathematics.

\begin{equation}
  \label{eq:ch01:k4:integral}
  \int_{0}^{1} x^{n}\,dx = \frac{1}{n+1}.
\end{equation}

The geometric series in Equation~\ref{eq:ch01:k4:series} converges for
$|r| < 1$ and gives a second display to render in this chapter.

\begin{equation}
  \label{eq:ch01:k4:series}
  \sum_{k=0}^{m} r^{k} = \frac{1-r^{m+1}}{1-r}.
\end{equation}

\subsubsection{A table}
\begin{table}[h]
  \centering
  \begin{tabular}{lrr}
    \toprule
    Run & Time & Pages \\
    \midrule
    One   & 1.20 & 3 \\
    Two   & 1.15 & 3 \\
    Three & 1.18 & 3 \\
    \bottomrule
  \end{tabular}
  \caption{A timing sample for chapter 1.}
  \label{tab:ch01:k4:runs}
\end{table}

Table~\ref{tab:ch01:k4:runs} is a valid float with a caption and a label. It
exists so the page builder has a boxed object that cannot be split across a
page break. The bibliography entry \cite{knuth1984} is resolved by the same compile.

\subsubsection{More prose}
A benchmark chapter needs enough text that the type\-set is visible above the
engine boot. This paragraph, the next one and the table above add up to a few
pages per chapter.

The content is intentionally repetitive across chapters. The point is not the
prose; it is that the include graph, the equation numbering and the table
labels all have to stay consistent while the probe edits one letter.

The final paragraph of chapter 1 gives the probe a second stable offset. A
change here must not invalidate the labels in any other chapter.

\subsection{Part 5}

This is included chapter 1 of twelve in
the large MDraft LuaLaTeX benchmark. Ordinary prose keeps the page count
realistic and gives the probe a stable word to edit. Equation~\ref{eq:ch12:k5:integral} and
Table~\ref{tab:ch12:k5:runs} are valid cross-chapter references.

\subsubsection{An integral}
The integral in Equation~\ref{eq:ch01:k5:integral} is valid for every integer
$n \ge 0$. It is deliberately simple; the benchmark measures the engine, not
the mathematics.

\begin{equation}
  \label{eq:ch01:k5:integral}
  \int_{0}^{1} x^{n}\,dx = \frac{1}{n+1}.
\end{equation}

The geometric series in Equation~\ref{eq:ch01:k5:series} converges for
$|r| < 1$ and gives a second display to render in this chapter.

\begin{equation}
  \label{eq:ch01:k5:series}
  \sum_{k=0}^{m} r^{k} = \frac{1-r^{m+1}}{1-r}.
\end{equation}

\subsubsection{A table}
\begin{table}[h]
  \centering
  \begin{tabular}{lrr}
    \toprule
    Run & Time & Pages \\
    \midrule
    One   & 1.20 & 3 \\
    Two   & 1.15 & 3 \\
    Three & 1.18 & 3 \\
    \bottomrule
  \end{tabular}
  \caption{A timing sample for chapter 1.}
  \label{tab:ch01:k5:runs}
\end{table}

Table~\ref{tab:ch01:k5:runs} is a valid float with a caption and a label. It
exists so the page builder has a boxed object that cannot be split across a
page break. The bibliography entry \cite{knuth1984} is resolved by the same compile.

\subsubsection{More prose}
A benchmark chapter needs enough text that the type\-set is visible above the
engine boot. This paragraph, the next one and the table above add up to a few
pages per chapter.

The content is intentionally repetitive across chapters. The point is not the
prose; it is that the include graph, the equation numbering and the table
labels all have to stay consistent while the probe edits one letter.

The final paragraph of chapter 1 gives the probe a second stable offset. A
change here must not invalidate the labels in any other chapter.

\subsection{Part 6}

This is included chapter 1 of twelve in
the large MDraft LuaLaTeX benchmark. Ordinary prose keeps the page count
realistic and gives the probe a stable word to edit. Equation~\ref{eq:ch12:k6:integral} and
Table~\ref{tab:ch12:k6:runs} are valid cross-chapter references.

\subsubsection{An integral}
The integral in Equation~\ref{eq:ch01:k6:integral} is valid for every integer
$n \ge 0$. It is deliberately simple; the benchmark measures the engine, not
the mathematics.

\begin{equation}
  \label{eq:ch01:k6:integral}
  \int_{0}^{1} x^{n}\,dx = \frac{1}{n+1}.
\end{equation}

The geometric series in Equation~\ref{eq:ch01:k6:series} converges for
$|r| < 1$ and gives a second display to render in this chapter.

\begin{equation}
  \label{eq:ch01:k6:series}
  \sum_{k=0}^{m} r^{k} = \frac{1-r^{m+1}}{1-r}.
\end{equation}

\subsubsection{A table}
\begin{table}[h]
  \centering
  \begin{tabular}{lrr}
    \toprule
    Run & Time & Pages \\
    \midrule
    One   & 1.20 & 3 \\
    Two   & 1.15 & 3 \\
    Three & 1.18 & 3 \\
    \bottomrule
  \end{tabular}
  \caption{A timing sample for chapter 1.}
  \label{tab:ch01:k6:runs}
\end{table}

Table~\ref{tab:ch01:k6:runs} is a valid float with a caption and a label. It
exists so the page builder has a boxed object that cannot be split across a
page break. The bibliography entry \cite{knuth1984} is resolved by the same compile.

\subsubsection{More prose}
A benchmark chapter needs enough text that the type\-set is visible above the
engine boot. This paragraph, the next one and the table above add up to a few
pages per chapter.

The content is intentionally repetitive across chapters. The point is not the
prose; it is that the include graph, the equation numbering and the table
labels all have to stay consistent while the probe edits one letter.

The final paragraph of chapter 1 gives the probe a second stable offset. A
change here must not invalidate the labels in any other chapter.

\subsection{Part 7}

This is included chapter 1 of twelve in
the large MDraft LuaLaTeX benchmark. Ordinary prose keeps the page count
realistic and gives the probe a stable word to edit. Equation~\ref{eq:ch12:k7:integral} and
Table~\ref{tab:ch12:k7:runs} are valid cross-chapter references.

\subsubsection{An integral}
The integral in Equation~\ref{eq:ch01:k7:integral} is valid for every integer
$n \ge 0$. It is deliberately simple; the benchmark measures the engine, not
the mathematics.

\begin{equation}
  \label{eq:ch01:k7:integral}
  \int_{0}^{1} x^{n}\,dx = \frac{1}{n+1}.
\end{equation}

The geometric series in Equation~\ref{eq:ch01:k7:series} converges for
$|r| < 1$ and gives a second display to render in this chapter.

\begin{equation}
  \label{eq:ch01:k7:series}
  \sum_{k=0}^{m} r^{k} = \frac{1-r^{m+1}}{1-r}.
\end{equation}

\subsubsection{A table}
\begin{table}[h]
  \centering
  \begin{tabular}{lrr}
    \toprule
    Run & Time & Pages \\
    \midrule
    One   & 1.20 & 3 \\
    Two   & 1.15 & 3 \\
    Three & 1.18 & 3 \\
    \bottomrule
  \end{tabular}
  \caption{A timing sample for chapter 1.}
  \label{tab:ch01:k7:runs}
\end{table}

Table~\ref{tab:ch01:k7:runs} is a valid float with a caption and a label. It
exists so the page builder has a boxed object that cannot be split across a
page break. The bibliography entry \cite{knuth1984} is resolved by the same compile.

\subsubsection{More prose}
A benchmark chapter needs enough text that the type\-set is visible above the
engine boot. This paragraph, the next one and the table above add up to a few
pages per chapter.

The content is intentionally repetitive across chapters. The point is not the
prose; it is that the include graph, the equation numbering and the table
labels all have to stay consistent while the probe edits one letter.

The final paragraph of chapter 1 gives the probe a second stable offset. A
change here must not invalidate the labels in any other chapter.

\subsection{Part 8}

This is included chapter 1 of twelve in
the large MDraft LuaLaTeX benchmark. Ordinary prose keeps the page count
realistic and gives the probe a stable word to edit. Equation~\ref{eq:ch12:k8:integral} and
Table~\ref{tab:ch12:k8:runs} are valid cross-chapter references.

\subsubsection{An integral}
The integral in Equation~\ref{eq:ch01:k8:integral} is valid for every integer
$n \ge 0$. It is deliberately simple; the benchmark measures the engine, not
the mathematics.

\begin{equation}
  \label{eq:ch01:k8:integral}
  \int_{0}^{1} x^{n}\,dx = \frac{1}{n+1}.
\end{equation}

The geometric series in Equation~\ref{eq:ch01:k8:series} converges for
$|r| < 1$ and gives a second display to render in this chapter.

\begin{equation}
  \label{eq:ch01:k8:series}
  \sum_{k=0}^{m} r^{k} = \frac{1-r^{m+1}}{1-r}.
\end{equation}

\subsubsection{A table}
\begin{table}[h]
  \centering
  \begin{tabular}{lrr}
    \toprule
    Run & Time & Pages \\
    \midrule
    One   & 1.20 & 3 \\
    Two   & 1.15 & 3 \\
    Three & 1.18 & 3 \\
    \bottomrule
  \end{tabular}
  \caption{A timing sample for chapter 1.}
  \label{tab:ch01:k8:runs}
\end{table}

Table~\ref{tab:ch01:k8:runs} is a valid float with a caption and a label. It
exists so the page builder has a boxed object that cannot be split across a
page break. The bibliography entry \cite{knuth1984} is resolved by the same compile.

\subsubsection{More prose}
A benchmark chapter needs enough text that the type\-set is visible above the
engine boot. This paragraph, the next one and the table above add up to a few
pages per chapter.

The content is intentionally repetitive across chapters. The point is not the
prose; it is that the include graph, the equation numbering and the table
labels all have to stay consistent while the probe edits one letter.

The final paragraph of chapter 1 gives the probe a second stable offset. A
change here must not invalidate the labels in any other chapter.
